\cotefile{18}
\scanpage{18}
A - 20 - \hfill n° 413

(resp. $U_n(\underline V)$) et ses éléments seront appelés \uline{distributions} (resp. \uline{distributions d'ordre} $\leq n$) sur $\underline V$. La restriction $\varphi_*$ à $U(\underline V)$ du transposé de $\varphi^*$ est donc un isomorphisme de modules filtrés de $U(\underline V)$ sur $S(M)$ et permet de transporter à $U(\underline V)$ la structure de coalgèbre avec forme unité de $S(M)$ : la prop. 6 montre que la structure de coalgèbre obtenue est indépendante de la carte choisie et la multiplication dans $\widehat P(\underline V)$ (qui s'identifie au dual de $U(\underline V)$) est la transposée de l'application diagonale de $U(\underline V)$. On désignera par $U_+(\underline V)$ le noyau de la forme unité ; c'est aussi l'image réciproque par $\varphi_*$ du sous-module $\displaystyle\sum_{k>0}S^k(M)$. On appellera \uline{module tangent} à $V$ et on notera $T(\underline V)$ l'intersection $U_+(\underline V)\cap U_1(\underline V)$ : c'est aussi l'image réciproque de $M=S^1(M)$ par $\varphi_*$.

Si $f\in\widehat P_+(\underline V;\underline V')$, le transposé de $f^*$ envoie $U(\underline V)$ dans $U(\underline V')$. Sa restriction à $U(\underline V)$ sera notée $f_*$. C'est une application linéaire compatible avec les filtrations, avec les structures de coalgèbres avec forme unité. La restriction de $f_*$ à $T(\underline V)$ est une application linéaire de $T(\underline V)$ dans $T(\underline V')$ appelée \uline{application linéaire tangente à} $f$ et notée $T(f)$ : si $\underline V$ et $\underline V'$ sont les noyaux formels définis par des modules $M$ et $M'$, cette définition est en accord avec celle donnée au n° 1, à condition d'identifier $T(M)$ et $M$ (et $T(M')$ avec $M'$) ce qui est loisible.

Supposons maintenant que $\underline V'$ soit le noyau formel associé à un module libre $M$. On peut alors à $\underline u\in\widehat P_+(\underline V;M)$ associer la composée de $\underline u_*$ et de la projection canonique de $S(M)$ sur $S^1(M)=M$ : on obtient ainsi un isomorphisme de $\widehat P_+(\underline V;M)$ sur le sous-module de $\operatorname{Hom}(U(\underline V);M)$ formé des applications linéaires nulles sur $U_0(\underline V)$ (il suffit pour le voir de prendre une carte de $\underline V$ et d'utiliser la remarque 2 du n° 4). Pour tout module séparé $N$, nous pouvons définir une application formelle

\margin{\mspage{\unsure{17}}}
